
\subsection{\texorpdfstring{\textbf{
                  Применение непрерывной сети Хопфилда к задаче LMAPF}}
      {Применение непрерывной сети Хопфилда к задаче LMAPF}}

Для применения полученных результатов взята задача LMAPF -- lifelong
multiagent pathfinding, где требуется одним из возможных критериев
хорошего решения является пропускная способность. То есть измеряется
отношение количества позиций, достигнутых агентами за определённый
период времени, к этому периоду времени. Чем оно больше, тем лучше
решена задача. Конечно, существуют различные факторы. Как, например,
плотность препятствий, правила начального расположения агентов,
непрерывный или дискретный случай в пространстве или на графе.

Здесь хочется упомянуть, что на начало исследования не было известно,
что подобная задача была решена для двумерной сетки с помощью
оптимизированной сети Хопфилда. Например, в работе Мерт Туранли
``Энергоэффективное управление покрытием с использованием нескольких
роботов на основе сетей Хопфилда'' \cite{turanli2020multi}

или в работе Юдинцева Б. С. ''Интеллектуальная система планирования
траекторий мобильных роботов, построенная на сети хопфилда'' \cite{yudintsev2014intelligent}

А также, в работе ''Обход препятствий с помощью модифицированной
нейронной сети Хопфилда'' \cite{kin2008obstacle}

Однако, в этой работе основной целью исследования является применимость
алгоритма скоростного градиента к обучению непрерывной модели Хопфилда.

В данной работе, рассматривается непрерывный вариант на ограниченном
двумерном поле с препятствиями в виде случайно расставленных
параллелепипедов. Каждый агент имеет радиус, ограниченные вращение и
скорость движения.

В качестве основы выбрано решение этой задачи, которое называют MPC --
Model Predictive Control, основывающееся на минимизации стоимости
функционала локального управления на каждом шаге симуляции.

Функционал стоимости состоит из нескольких компонент, которые
суммируются:

\begin{enumerate}
      \def\labelenumi{\arabic{enumi}.}
      \item
            J. Ошибка цели. Близость агента к цели.
      \item
            H. Ошибка блокировки. Агент может вынужденно обойти препятствие, отдалившись от цели.
      \item
            G. Ошибка близости к препятствиям. Агент должен держаться подальше от препятствий.
      \item
            D. Ошибка близости к агентам. Агент должен держаться подальше от агентов.
\end{enumerate}

Вычисление задачи минимизации будет решением задачи локального
управления. \cite{ota2022proposal}

Оно дополняется глобальным поиском путей, чтобы агент не
попадал в тупики без выходов и хорошо ориентировался в больших
пространствах. Как дополнение, может появиться алгоритм, который будет
управлять фазовыми переходами всей системы, посылая сигналы конкретным
агентам об изменении глобального пути.

Использовать нейронную сеть Хопфилда планировалось как дополнение -
компоненту функционала стоимости -- локального выбора пути и обучать
прямо в ходе выполнения симуляции.

Конкретнее, в локальных координатах агента (с его вращением) строится
двумерная сетка. Каждая её ячейка получает свой вес -- цену прохода,
которая вычисляется также на каждом шаге симуляции. Цена вычисляется
исходя из препятствий, агентов, а также памяти агента. Веса же
становятся частью сети Хопфилда, которая находит путь наименьшей
стоимости, предсказывая будущие положения агентов.
